%TCIDATA{Version=5.50.0.2890}
%TCIDATA{LaTeXparent=0,0,NoticeDKE.tex}


Basically, the determination of the steady-state electron distribution
function is an initial value problem, where the initial guess corresponds
usually to the unperturbed solution of the linearized problem. For the
zero-order Fokker-Planck equation, the Maxwellian distribution function is
therefore by definition an eigenfunction of the collision operator, without
external perturbation. Here the weakly relativistic solution corresponding
to the condition on the relativistic factor $\gamma -1\ll 1$ is considered
at time $t=0$,
\begin{equation}
f_{M}\left( \psi ,p\right) \simeq \frac{n_{e}\left( \psi \right) }{\left[ 2%
\pi T_{e}\left( \psi \right) \right] ^{3/2}}\exp \left[ -\frac{p^{2}}{\left(
1+\gamma \right) T_{e}\left( \psi \right) }\right]
\end{equation}
where $T_{e}\left( \psi \right) $ and $n_{e}\left( \psi \right) $ are the
electron temperature and density respectively. Details on the notation are
given in Sec. \ref{normtene}. Projected on the numerical grids, as defined
in Sec.\ref{Griddef}, the discrete form is
\begin{equation}
f_{M,l+1/2,i+1/2,j+1/2}\simeq \frac{n_{e,l+1/2}}{\left[ 2\pi %
T_{e,l+1/2}\right] ^{3/2}}\exp \left[ -\frac{p_{i+1/2}^{2}}{\left( 1+\gamma
_{i+1/2}\right) T_{e,l+1/2}}\right]
\end{equation}

When the Lower Hybrid current drive problem is addressed, it is possible to
start from a guess that already incorporate the existence of a plateau
region, using the formulation given in Ref. \cite{fuc85} in the
non-relativistic limit. When relativistic corrections must be considered,
the effective perpendicular temperature $T_{\bot }$ in the resonance domain
must be usualy  multipled by $2$. Though this elegant approach seems
attractive in order to reduce the total number of iterations for reaching
the steady-state solution, its effectiveness is usualy poor when large time
step $\Delta t$ are considered. Indeed, the total number of time steps is
merely determined by the relaxation time of the most energetic part of the
fast electron tail prodiced by the wave. Since electrons are very weakly
collisional, their relaxation time is very long as compared to the thermal
one. Therefore, the gain for the rate of convergence is usually very small
since the distribution model does not describe accurately the region above
the plateau of the momentum space. Consequently, even if this method is
implemented in the code, it is almost never used.
